\documentclass{article}

\usepackage{float}
\usepackage{arxiv}

\usepackage[utf8]{inputenc}
\usepackage[T1]{fontenc}
\usepackage{hyperref}
\usepackage{url}
\usepackage{booktabs}
\usepackage{amsfonts}
\usepackage{nicefrac}
\usepackage{microtype}
\usepackage{graphicx}
\usepackage{listings}
\usepackage{xcolor}
\usepackage{fancyhdr}
\graphicspath{ {./neatagent_images/} }

\lstset{
  basicstyle=\ttfamily\small,
  backgroundcolor=\color{gray!10},
  frame=single,
  breaklines=true,
  columns=fullflexible,
}

\pagestyle{fancy}
\fancyhf{}
\rfoot{\thepage}
\renewcommand{\headrulewidth}{0pt}

\title{NeatAgent: A Local AI Agent for Autonomous File Organization}

\author{
  Abdulmohsen Alghannam \\
  Columbia University \\
  \texttt{afa2165@columbia.edu} \\
}

\begin{document}
\maketitle
\thispagestyle{fancy}

% ─── Synopsis ────────────────────────────────────────────────────────────────
\section*{Synopsis}

NeatAgent is a local AI agent that organizes files on your machine. You point it at a folder, it reads the contents of your files, groups them into meaningful subfolders, renames anything with a generic name, shows you the plan, and executes it after you approve. Nothing leaves your machine — no cloud API, no subscription, no internet connection needed.

The agent is built on a \textbf{PLAN-AND-ACT} loop: it creates a complete plan before touching anything, then executes one step at a time. If a step fails, a Reflector node decides whether to retry, skip, or replan. The whole thing is wired together using \textbf{LangGraph}, a library for building graph-based agent workflows, and runs inference locally through \textbf{Ollama}, which serves open-weight language models on your own hardware.

One thing that made the project interesting was figuring out how to get small local models to produce reliable, structured output. I went through a few iterations before landing on the current design, and a lot of the work ended up being about preventing the model from hallucinating files that don't exist. I'll explain how I solved that in the Novelty section.

For evaluation, I created a dataset of 10 files with generic names (file1.txt, file4.pdf, etc.) covering different content types — invoices, health reports, gym logs, photos, and CSVs. The goal was to see whether the agent could classify files correctly based on \textit{content} alone, since the filenames give no signal. I compared three models of different sizes: \texttt{qwen2.5-coder:14b}, \texttt{qwen3:8b}, and \texttt{llama3.2:3b}.

\begin{figure}[H]
    \centering
    \includegraphics[width=0.85\linewidth]{NeatAgent_architecture.png}
    \caption{NeatAgent architecture — Planner, Executor, and Reflector nodes connected through a shared state}
    \label{fig:architecture}
\end{figure}

% ─── Section 1: System Design ────────────────────────────────────────────────
\section{System Design}

The agent has three main nodes that communicate through a shared state dictionary (\texttt{AgentState}):

\begin{itemize}
    \item \textbf{Planner} — lists the files in the target folder, reads content previews for files with ambiguous names, and asks the LLM for a \texttt{\{filename $\to$ folder\}} mapping. Python code converts that mapping into a step list.
    \item \textbf{Executor} — runs one step per graph invocation (create folder, move file, rename file). Supports dry-run mode (shows plan only) and safe mode (asks before each destructive step).
    \item \textbf{Reflector} — only fires when a step fails. It auto-skips errors that will never recover (e.g., file not found), and for everything else asks the LLM whether to retry, skip, or replan.
\end{itemize}

The graph loops the Executor via a conditional edge — not a Python for-loop — so the Reflector can intercept failures between steps. A 400-second wall-clock timeout prevents the agent from running indefinitely on large folders.

\begin{figure}[H]
    \centering
    \includegraphics[width=0.85\linewidth]{NeatAgent_plan.png}
    \caption{The plan presented to the user before execution. The user can approve, reject, or switch to dry-run mode.}
    \label{fig:plan}
\end{figure}

\subsection{What I Built On}

\begin{table}[H]
\centering
\caption{External Libraries and Frameworks}
\begin{tabular}{lll}
\toprule
\textbf{Component} & \textbf{Library} & \textbf{What it does in NeatAgent} \\
\midrule
Agent framework     & LangGraph \cite{langgraph}       & Wires the Planner/Executor/Reflector nodes together \\
LLM serving         & Ollama                           & Runs quantized LLMs locally via REST API \\
LLM client          & langchain-ollama                 & \texttt{ChatOllama} used in Planner and Reflector \\
PDF reading         & pdfplumber                       & Extracts text from PDFs for content-based classification \\
DOCX reading        & python-docx                      & Extracts text from Word documents \\
Terminal UI         & rich                             & Colored plan table, progress display, summary \\
Interactive prompts & questionary                      & Folder picker and run-option checkboxes \\
\midrule
\multicolumn{3}{l}{\textit{Models used}} \\
\midrule
Primary     & \texttt{qwen2.5-coder:14b} (Alibaba) & Best JSON output reliability \\
Mid-size    & \texttt{qwen3:8b} (Alibaba)          & Faster, slightly lower accuracy \\
Smallest    & \texttt{llama3.2:3b} (Meta)          & Fastest, lowest accuracy \\
Vision      & \texttt{qwen3-vl:8b} (Alibaba)       & Describes images for classification \\
\bottomrule
\end{tabular}
\end{table}

The PLAN-AND-ACT pattern is drawn from Erdogan et al.\ \cite{erdogan2025planact}. The Reflector is modeled after the Reflexion framework \cite{anthropic2024agents}, where an agent reflects on its own failures before deciding what to do next.

% ─── Section 2: Novelty ──────────────────────────────────────────────────────
\section{What's Novel}

\subsection{Mapping-Based Planner}

The first version of the planner asked the LLM to generate a full list of action steps: move this file here, rename that one, etc. This turned out to be a disaster — on a 10-file folder the model produced 63 steps, 40 of which failed because the model had invented filenames like \texttt{file11.pdf} through \texttt{file52.kt} that didn't exist. The model would drift after covering the real files and keep generating steps for imaginary ones.

The fix was to change what the LLM is responsible for. Instead of generating steps, it now outputs a simple mapping:

\begin{lstlisting}[language={}]
{
  "file1.txt": { "folder": "Invoices", "rename": null },
  "file4.pdf": { "folder": "Invoices", "rename": null },
  "file6.jpg": { "folder": "Photos",   "rename": null }
}
\end{lstlisting}

The keys come from the real file listing, so the model can't invent a key that doesn't exist in its input. Python then generates all the steps deterministically. Any key that doesn't match a real file is filtered out before any steps are built. This made hallucination structurally impossible rather than just something we hoped the model would avoid.

\subsection{Content-Aware Classification}

For files with generic names, the filename alone gives no signal. \texttt{file4.pdf} could be anything. NeatAgent solves this by reading a content preview before calling the planner: text files and PDFs are read with pdfplumber, DOCX files with python-docx, and images are described by a vision model. The preview is injected into the planner prompt as a \texttt{CONTENT:} hint. This gave a roughly 20 percentage point improvement in placement accuracy for ambiguous files.

\subsection{Fully Local, Including the Evaluator}

Everything runs locally — not just the agent, but also the LLM-as-judge evaluator. The judge is just another Ollama model that receives the resulting folder structure and ground truth, and returns a 1–10 score with reasoning. No API key needed at any point.

\begin{figure}[H]
    \centering
    \includegraphics[width=0.85\linewidth]{NeatAgent_CLI.png}
    \caption{The interactive setup wizard launched on first run via \texttt{python run.py}}
    \label{fig:cli}
\end{figure}

% ─── Section 3: Research Questions ───────────────────────────────────────────
\section{Research Questions}

\subsection{RQ1: Can local LLMs organize files without cloud APIs?}

\textbf{Motivation.} Most AI agent systems require a cloud API (GPT-4, Claude). The question is whether a quantized local model can do this well enough to be useful.

\textbf{Results.} Using \texttt{qwen2.5-coder:14b} on the 10-file \texttt{model\_test} dataset, the agent achieved 100\% coverage (every file was moved) and around 60\% placement accuracy (6 out of 10 files in the correct semantic folder). The agent correctly grouped invoices, photos, and CSVs. The main misclassifications involved files where the content could reasonably belong in more than one category.

\textbf{Answer.} Yes. A 14B local model gets the job done without cloud dependency and produces results a typical user would find useful.

\subsection{RQ2 (Ablation): Which parts of the pipeline actually matter?}
\label{sec:ablation}

I wanted to understand which components were doing the real work versus which ones were just overhead. I tested four conditions:

\begin{enumerate}
    \item \textbf{Step-based planner (original)} — LLM generates executable steps directly, no content preview, no reflector
    \item \textbf{Mapping planner, no content preview} — LLM sees filenames and sizes only
    \item \textbf{Mapping planner + content preview, no reflector} — reflector disabled, failed steps auto-skipped
    \item \textbf{Full system} — mapping planner + content preview + reflector
\end{enumerate}

\begin{table}[H]
\centering
\caption{Ablation Results — \texttt{model\_test} dataset, \texttt{qwen2.5-coder:14b}}
\begin{tabular}{lrrrr}
\toprule
\textbf{Condition} & \textbf{Steps} & \textbf{Failures} & \textbf{Hallucinated Files} & \textbf{Placement Acc.} \\
\midrule
(1) Step-based (baseline) & 63 & 40 & $\sim$42 & unusable \\
(2) Mapping, no preview   & 13 &  0 & 0        & $\sim$40\% \\
(3) Mapping + preview, no reflector & 13 & 0 & 0 & $\sim$60\% \\
(4) Full system           & 13 &  0 & 0        & $\sim$60\% \\
\bottomrule
\end{tabular}
\end{table}

The mapping-based planner was the biggest single change — it took the system from completely broken (condition 1) to functional (condition 2). Content previews added about 20 points of placement accuracy for files with ambiguous names. The Reflector didn't change the numbers on the clean \texttt{model\_test} runs, but it recovered gracefully during development when files were already moved or folders already existed.

\subsection{RQ3: Does model size make a meaningful difference?}

\begin{table}[H]
\centering
\caption{Three models on \texttt{model\_test} — full pipeline}
\begin{tabular}{lrrrr}
\toprule
\textbf{Model} & \textbf{Size} & \textbf{Coverage} & \textbf{Placement Acc.} & \textbf{Time} \\
\midrule
\texttt{qwen2.5-coder:14b} & 14B & 100\% & $\sim$60\% & $\sim$90s \\
\texttt{qwen3:8b}          &  8B & 100\% & $\sim$50\% & $\sim$70s \\
\texttt{llama3.2:3b}       &  3B & 100\% & $\sim$40\% & $\sim$40s \\
\bottomrule
\end{tabular}
\end{table}

All three models moved every file (100\% coverage). Accuracy scaled with size, but even the 3B model produced a reasonable organization — it just tended to use broader folder names and occasionally confused related categories (e.g., putting a health CSV into a Fitness folder instead of Data).

\begin{figure}[H]
    \centering
    \includegraphics[width=0.85\linewidth]{NeatAgent_3models_compared.png}
    \caption{Placement accuracy and inference time across three model sizes}
    \label{fig:models}
\end{figure}

% ─── Section 4: Value to Users ───────────────────────────────────────────────
\section{Value to Users}

NeatAgent is for anyone who accumulates disorganized files and doesn't want to manually sort them. The privacy angle is important: if you work with sensitive documents — contracts, medical files, financial records — you probably don't want to upload them to a cloud API to classify them. NeatAgent keeps everything on your machine.

The setup is designed to be as friction-free as possible. Running \texttt{python run.py} handles virtual environment creation, dependency installation, and first-time Ollama setup automatically. After that, you pick a folder from a menu and the agent does the rest.

\textbf{Availability:}
\begin{itemize}
    \item Open-source: \url{https://github.com/Abdulmohseng/6156-Agent-Project}
    \item Works on macOS, Linux, and Windows (Python 3.10+)
    \item Free — no API keys, no subscriptions
    \item Works offline after models are downloaded
    \item Hardware: 8–16 GB RAM is enough for the 3B–8B models; 14B needs about 10 GB
\end{itemize}

% ─── Section 5: Evaluation and Metrics ───────────────────────────────────────
\section{Evaluation and Metrics}

\subsection{The Dataset I Created}

I built a 10-file dataset (\texttt{model\_test}) specifically for this evaluation. The goal was to make classification hard: every file has a generic name (file1.txt through file10.csv), so the agent \textit{has} to read the content to know where to put it. The files cover five categories:

\begin{itemize}
    \item \textbf{Invoices} — \texttt{file1.txt} (text invoice), \texttt{file4.pdf} (PDF invoice)
    \item \textbf{Health} — \texttt{file3.txt} (gym log), \texttt{file5.pdf} (annual health report)
    \item \textbf{Personal} — \texttt{file2.txt} (book reading log)
    \item \textbf{Photos} — \texttt{file6.jpg}, \texttt{file7.jpg}, \texttt{file8.jpg}
    \item \textbf{Data} — \texttt{file9.csv} (sensor readings), \texttt{file10.csv} (sales data)
\end{itemize}

Each file has a ground truth label with a primary expected folder and a list of aliases — folder names the agent might reasonably use that should still count as correct. For example, an invoice file accepts ``Invoices'', ``Finance'', ``Bills'', and ``billing'' as all correct. I also have a larger dataset (\texttt{sample\_flat}, 24 files) with a mix of generic and descriptive filenames.

\subsection{Evaluation Pipeline}

\begin{figure}[H]
    \centering
    \includegraphics[width=0.85\linewidth]{NeatAgent_evaluation.png}
    \caption{Evaluation pipeline output — per-file placement results and aggregate scores}
    \label{fig:eval}
\end{figure}

The pipeline has four parts:

\begin{enumerate}
    \item \textbf{\texttt{runner.py}} — runs the agent as a subprocess with auto-approval, then finds the manifest it wrote to \texttt{\textasciitilde/.file-agent/runs/}
    \item \textbf{\texttt{compare.py}} — walks the execution trace to find where each file ended up, including through rename chains
    \item \textbf{\texttt{judge.py}} — sends the result to a local Ollama model and asks for a 1–10 score with reasoning
    \item \textbf{\texttt{report.py}} — prints a colored terminal table and saves JSON
\end{enumerate}

\subsection{Metrics}

\textbf{Coverage} — what fraction of files were moved at all:
\[\text{Coverage} = \frac{|\{f : \text{final\_folder}(f) \neq \text{root}\}|}{|\text{files}|}\]

\textbf{Placement Accuracy} — what fraction landed in the right folder (exact or alias match):
\[\text{Accuracy} = \frac{|\{f : \text{final\_folder}(f) \in \text{aliases}(f)\}|}{|\text{files}|}\]

\textbf{LLM-as-Judge Score (1--10)} — a local Ollama model sees the organized folder structure and the ground truth, and scores the overall quality with reasoning. Scores of 4--5/10 came back consistently, which was lower than the structural metrics suggested. After looking at the reasoning, the judge was penalizing folder names like ``Invoices'' when it would have preferred ``Finance'' — even though both are correct. This highlights that file organization is inherently subjective.

% ─── Section 6: Deliverables ─────────────────────────────────────────────────
\section{Deliverables}

\textbf{Repository:} \url{https://github.com/Abdulmohseng/6156-Agent-Project}

\begin{table}[H]
\centering
\caption{Project Deliverables}
\begin{tabular}{ll}
\toprule
\textbf{Deliverable} & \textbf{Location} \\
\midrule
Project proposal          & \texttt{docs/project\_proposal.pdf} \\
Progress report           & \texttt{docs/NeatAgent.pdf} \\
Final report              & \texttt{docs/final\_report.md} \\
This paper                & \texttt{docs/paper.tex} \\
Agent source code         & \texttt{file-agent/} \\
Entry point               & \texttt{run.py} \\
First-time setup wizard   & \texttt{setup/} \\
Evaluation pipeline       & \texttt{eval/} \\
Ground truth labels       & \texttt{eval/ground\_truth/} \\
Test datasets             & \texttt{tests/data/model\_test/}, \texttt{eval/datasets/sample\_flat/} \\
Evaluation run results    & \texttt{eval/results/} \\
Build \& run instructions & \texttt{README.md} \\
\bottomrule
\end{tabular}
\end{table}

% ─── Section 7: Self-Evaluation ──────────────────────────────────────────────
\section{Self-Evaluation}

This was an individual project. I built the full system from scratch: the LangGraph state machine, the planner rewrite (twice — once step-based, once mapping-based), the executor with dry-run and safe mode, the reflector, five filesystem tools, the evaluation pipeline, two labeled datasets, the setup wizard, and the documentation.

\subsection{What I Found Hard}

The hallucination problem was the biggest challenge and took the longest to figure out. The step-based planner seemed to work fine on simple test cases during development, so I only discovered the issue when I ran formal evaluation on a slightly more complex folder. At that point the model was producing 63 steps for 10 files and failing 40 of them. Tracking down why — reading the raw LLM output, understanding the context drift, realizing the model was continuing to generate after it ran out of real files — took a while. The fix (changing the output format) was simple once I understood the problem, but getting there wasn't.

Evaluation design was also harder than I expected. Defining ``correct'' for a subjective task like file organization requires a lot of judgment calls. Building the alias system helped, but it's still imperfect — a real evaluation would benefit from multiple annotators comparing notes.

\subsection{What I Learned}

The most lasting thing I took from this project is that output format is more important than model capability for agentic tasks. I spent time trying to fix the hallucination problem through better prompting — more explicit instructions, stronger warnings, different phrasing. None of it helped. Changing the \textit{structure} of what the model was asked to output fixed it immediately. When you can constrain the output schema so that the error can't be represented, you've solved the problem in a way that's reliable regardless of which model you're using.

The other thing I learned is that small local models are more capable than I expected for tasks with a well-defined output format. A 3B model running on a laptop can organize files semantically. That's genuinely useful and opens up automation for people who can't or won't use cloud APIs.

% ─── Bibliography ────────────────────────────────────────────────────────────
\bibliographystyle{unsrt}

\begin{thebibliography}{6}

\bibitem{erdogan2025planact}
Lutfi Eren Erdogan et al.
\newblock {PLAN-AND-ACT}: Improving Planning of Agents for Long-Horizon Tasks.
\newblock {\em arXiv preprint arXiv:2503.09572}, 2025.
\newblock \url{https://arxiv.org/pdf/2503.09572}

\bibitem{anthropic2024agents}
Anthropic.
\newblock Building Effective Agents.
\newblock Blog, 2024.
\newblock \url{https://www.anthropic.com/engineering/building-effective-agents}

\bibitem{langgraph}
LangChain.
\newblock LangGraph: Build Stateful, Multi-Actor Applications with LLMs.
\newblock \url{https://github.com/langchain-ai/langgraph}, 2024.

\bibitem{jin2024multiagent}
Taicheng Guo et al.
\newblock Large Language Model based Multi-Agents: A Survey of Progress and Challenges.
\newblock {\em arXiv preprint arXiv:2402.01680}, 2024.
\newblock \url{https://arxiv.org/pdf/2402.01680}

\bibitem{yang2024swe}
John Yang et al.
\newblock {SWE-agent}: Agent-Computer Interfaces Enable Automated Software Engineering.
\newblock In {\em NeurIPS}, 2024.
\newblock \url{https://arxiv.org/abs/2405.15793}

\bibitem{asai2023selfrag}
Akari Asai et al.
\newblock Self-{RAG}: Learning to Retrieve, Generate, and Critique through Self-Reflection.
\newblock {\em arXiv preprint arXiv:2310.11511}, 2023.
\newblock \url{https://arxiv.org/abs/2310.11511}

\end{thebibliography}

% ─── Appendix ────────────────────────────────────────────────────────────────
\clearpage
\appendix
\section*{Appendix}
\addcontentsline{toc}{section}{Appendix}

\subsection{Planner Prompt}

\begin{lstlisting}[language={}]
You are a file organization assistant. Given a list of files,
output a JSON object that maps each filename to its destination.

Output schema -- one key per file, exactly:
{
  "<exact_filename_from_listing>": {
    "folder": "<SemanticFolderName>",
    "rename": "<new_stem_no_extension_or_null>"
  }
}

Rules:
- Output a key for EVERY file. No more, no less. Never invent filenames.
- folder: short, meaningful name (Finance, Health, Work, Photos, etc.)
- rename: snake_case stem only -- set only when filename is generic.
- Use CONTENT and IMAGE hints to classify accurately.
Output ONLY valid JSON. No markdown, no explanation.
\end{lstlisting}

\subsection{Ground Truth Schema}

\begin{lstlisting}[language={}]
"file4.pdf": {
  "expected_folder": "invoices",
  "folder_aliases": ["Finance", "Bills", "Receipts"],
  "expected_name": null,
  "content_hint": "Invoice from Acme Corp, March 2026"
}
\end{lstlisting}

\subsection{Course Feedback}

The project structure was effective. Actually building an agent — not just reading about them — made it clear where the real difficulties are. One suggestion: a mid-semester peer review session would help students catch architectural mistakes before they're too embedded to change.

\end{document}
